
\subsection{Auditoria exploratória das coletas}
\label{subsec:auditoria_coletas_loso}

Além do filtro de pressão descrito acima, examinou-se a estabilidade dos
seis sinais MQ em cada coleta. Para o sensor $s$, a coleta $c$ e a leitura
$t>1$, calcularam-se a variação absoluta entre leituras consecutivas e o
limiar de cada sensor conforme a Equação~\eqref{eq:salto_mq_loso}:
\begin{equation}
    d_{s,c,t}=\left|x_{s,c,t}-x_{s,c,t-1}\right|,
    \qquad
    q_s=Q_{0{,}995}\!\left(\{d_{s,c,t}:c,\ t>1\}\right),
    \label{eq:salto_mq_loso}
\end{equation}
em que $x_{s,c,t}$ é o sinal bruto do MQ e $Q_{0{,}995}$ é o percentil
99,5 das variações reunidas de todas as coletas para aquele sensor.
As diferenças foram calculadas dentro de cada coleta, sem conectar o fim
de uma coleta ao início de outra. Definiu-se uma leitura com salto
multissensor quando pelo menos dois dos seis MQ superaram seus respectivos
limiares. O indicador de salto e a fração de leituras sinalizadas na coleta
foram calculados pela Equação~\eqref{eq:ruido_multissensor_loso}:
\begin{equation}
    J_{c,t}=\mathbf{1}\!\left\{
      \sum_{s\in\mathcal{S}_{\mathrm{MQ}}}
      \mathbf{1}\{d_{s,c,t}\geq q_s\}\geq 2
    \right\},
    \qquad
    R_c=\frac{1}{n_c}\sum_{t=1}^{n_c}J_{c,t},
    \label{eq:ruido_multissensor_loso}
\end{equation}
em que $n_c$ é o número de leituras da coleta. A regra exploratória de
qualidade sinalizou coletas com $R_c\geq 0{,}05$. A avaliação visual das
séries foi usada para contextualizar os saltos. O limiar foi aplicado à
coleta individual, de forma que as demais coletas do mesmo vaso permaneceram
elegíveis. Os percentis e a regra de exclusão foram definidos após a
inspeção da base completa; sua utilização constitui análise de
sensibilidade, e não um protocolo de qualidade estabelecido antes da
aquisição.

Separadamente, as previsões da análise inicial de Extra Trees com MQ2,
MQ3, MQ8, MQ135, MQ138 e as três entradas ambientais foram agrupadas por
coleta. A fração de erros de cada coleta foi calculada pela
Equação~\eqref{eq:erro_coleta_loso}:
\begin{equation}
    E_c=\frac{1}{n_c}\sum_{t=1}^{n_c}
    \mathbf{1}\{\widehat y_{c,t}\neq y_c\},
    \label{eq:erro_coleta_loso}
\end{equation}
em que $y_c$ é a classe experimental da coleta. Foram examinadas as
coletas com $E_c>0{,}5$. Na análise de sensibilidade com exclusões,
retiraram-se as quatro coletas identificadas por esse critério e uma
coleta adicional sinalizada pela auditoria multissensor; uma das quatro
também satisfazia o critério de ruído. A seleção por $E_c$ utilizou os
erros do próprio classificador na mesma base, e não demonstra defeito de
aquisição nas coletas selecionadas. Por isso, todas as comparações após
essas exclusões são exploratórias e condicionadas a essa seleção.

\subsection{Validação complementar com coletas agrupadas por vaso}
\label{subsec:loso_complementar}

Para examinar o efeito das coletas repetidas de um mesmo vaso na avaliação,
primeiro foram retiradas somente as cinco coletas selecionadas na análise
de sensibilidade. Nenhum vaso foi excluído: suas outras coletas
permaneceram na base. Em seguida, aplicou-se a validação
\textit{leave-one-vase-out} (LOSO) às coletas restantes, agrupadas pelo
vaso de origem. Em cada rodada, todas as coletas restantes de um mesmo
vaso foram separadas juntas para teste. O agrupamento utilizou a
combinação entre
classe experimental e código de vaso. Essa chave técnica evitou uma colisão
de identificadores na
planilha derivada, que contém registros da classe inoculada com código
numérico de vaso 1; no delineamento físico, os vasos 1 a 5 são controles
e os vasos 6 a 10 são inoculados. Para cada vaso $v$, a separação entre
as coletas de treinamento e de teste seguiu a
Equação~\eqref{eq:particao_loso}:
\begin{equation}
    \begin{aligned}
        \mathcal{C}_v
        &=\{c\in\mathcal{C}:v(c)=v\},\\
        \mathcal{D}^{\mathrm{teste}}_v
        &=\bigcup_{c\in\mathcal{C}_v}\mathcal{D}_c,\\
        \mathcal{D}^{\mathrm{treino}}_v
        &=\bigcup_{c\in\mathcal{C}\setminus\mathcal{C}_v}\mathcal{D}_c,
    \end{aligned}
    \label{eq:particao_loso}
\end{equation}
em que $\mathcal{C}$ é o conjunto de coletas elegíveis, $v(c)$ identifica
o vaso da coleta $c$ e $\mathcal{D}_c$ reúne suas leituras e classes.
Em cada uma das dez rodadas, treinou-se um modelo com as coletas dos
nove outros vasos e testou-se com todas as coletas elegíveis do vaso
reservado. Assim, coletas do mesmo vaso não participaram simultaneamente
do treinamento e do teste em uma mesma rodada, e cada leitura elegível
recebeu uma previsão fora do grupo usado no ajuste. O procedimento
preservou os dez vasos; apenas as cinco coletas retiradas antes do LOSO
deixaram de contribuir com leituras.

Para Extra Trees, usaram-se sinais MQ brutos, sem a compensação estatística
da Equação~\eqref{eq:compensacao_mq}, e, quando previstos na configuração,
o índice operacional do solo, a temperatura e a pressão. O classificador
foi reajustado em cada rodada com 900 árvores,
\texttt{max\_features="sqrt"}, \texttt{min\_samples\_leaf=10},
\texttt{bootstrap=False} e semente 42. A previsão por leitura foi a classe
com probabilidade estimada de classe 1 igual ou superior a 0,5.

Como comparação complementar, \textit{Random Forest} e a rede temporal
descrita anteriormente foram reajustados em cada rodada LOSO em duas
configurações fixas: MQ2, MQ3, MQ8, MQ135 e MQ138, isoladamente ou
acompanhados das três variáveis ambientais. O Random Forest usou 300
árvores, \texttt{class\_weight="balanced"} e semente 42. Na rede temporal,
a padronização das entradas foi estimada apenas no treinamento de cada
rodada. As janelas mantiveram 128 leituras e passo de 64, com cinco ou oito
canais conforme a configuração; nenhuma janela atravessou o limite de
uma coleta. Para uma leitura $i$ coberta pelo conjunto $\mathcal{W}(i)$
de janelas, a probabilidade da classe 1 e a previsão final foram obtidas
pela Equação~\eqref{eq:agregacao_janelas_loso}:
\begin{equation}
    \widehat p_i(1)=
    \frac{1}{|\mathcal{W}(i)|}
    \sum_{w\in\mathcal{W}(i)}\widehat p_w(1),
    \qquad
    \widehat y_i=\mathbf{1}\{\widehat p_i(1)\geq 0{,}5\}.
    \label{eq:agregacao_janelas_loso}
\end{equation}
Os modelos foram comparados por leitura após reunir as previsões de todas
as rodadas. Essa avaliação não corresponde à partição por coleta de 70\%/30\%
descrita anteriormente. A implementação \texttt{ExtraTreesClassifier}
descrita para a análise anterior é tecnicamente Extra Trees; nesta
comparação, \textit{Random Forest} designa especificamente a implementação
\texttt{RandomForestClassifier}. A rede temporal é identificada aqui como
TCN, pois a implementação usada não possui módulo de atenção ou de
explicabilidade integrado.

\subsection{Busca combinatória das entradas}
\label{subsec:busca_combinatoria_loso}

A busca de entradas foi realizada com Extra Trees e a mesma divisão LOSO.
As nove variáveis candidatas estão definidas na
Equação~\eqref{eq:conjunto_entradas_loso}:
\begin{equation}
    \mathcal{F}=\{\mathrm{MQ2},\mathrm{MQ3},\mathrm{MQ7},\mathrm{MQ8},
    \mathrm{MQ135},\mathrm{MQ138},I_{\mathrm{solo}},T,P\}.
    \label{eq:conjunto_entradas_loso}
\end{equation}
Foram avaliados todos os subconjuntos não vazios $A\subseteq\mathcal{F}$.
Cada uma das nove variáveis pode estar presente ou ausente, produzindo
$2^9$ subconjuntos. Um deles é o conjunto vazio, sem qualquer variável
de entrada para o classificador; por isso, ele foi excluído da busca.
O número de configurações avaliadas foi calculado pela
Equação~\eqref{eq:combinacoes_loso}:
\begin{equation}
    \left|\{A\subseteq\mathcal{F}:A\neq\emptyset\}\right|
    =2^{|\mathcal{F}|}-1=2^9-1=511.
    \label{eq:combinacoes_loso}
\end{equation}
Solo, temperatura e pressão variaram separadamente nessa busca,
inclusive nas configurações sem sensores MQ. O BMP280 fornece
temperatura e pressão como duas variáveis, embora seja um único componente
físico; portanto, a contagem de configurações acima se refere a entradas
do modelo, não a componentes de hardware. Se temperatura e pressão forem
tratadas como um bloco indivisível, passam a existir oito componentes
selecionáveis: seis sensores MQ, o sensor de solo e o BMP280. Cada
componente pode estar presente ou ausente, gerando $2^8$ combinações.
Excluindo a combinação vazia, sem nenhum componente, restam
$2^8-1=255$ configurações de componentes possíveis.

Para cada subconjunto $A$, o mesmo classificador e as mesmas dez partições
por vaso foram usados. A ordenação adotou a acurácia por leitura reunida
entre as rodadas; empates foram resolvidos por menor número total de
variáveis, depois por maior número de MQ e, por fim, por ordem alfabética
dos nomes das entradas. Random Forest e rede temporal não participaram
dessa busca exaustiva. A escolha da configuração e sua avaliação ocorreram
nos mesmos vasos, sem validação aninhada nem conjunto biológico externo;
portanto, a ordenação das configurações é exploratória.

\subsection{Cálculo das métricas na validação por vaso}
\label{subsec:metricas_loso_complementar}

Com as previsões das rodadas LOSO reunidas, a acurácia por leitura foi
calculada conforme a Equação~\eqref{eq:acuracia_loso}:
\begin{equation}
    \mathrm{Acc}=\frac{1}{N}
    \sum_{i=1}^{N}\mathbf{1}\{\widehat y_i=y_i\},
    \label{eq:acuracia_loso}
\end{equation}
em que $N$ é o total de leituras avaliadas. Para cada classe
$k\in\{0,1\}$, o recall, a precisão e a medida F1 foram calculados
pela Equação~\eqref{eq:metricas_classe_loso}:
\begin{equation}
    \mathrm{Recall}_k=\frac{TP_k}{TP_k+FN_k},\qquad
    \mathrm{Prec}_k=\frac{TP_k}{TP_k+FP_k},\qquad
    F1_k=\frac{2\,\mathrm{Prec}_k\,\mathrm{Recall}_k}
    {\mathrm{Prec}_k+\mathrm{Recall}_k},
    \label{eq:metricas_classe_loso}
\end{equation}
em que $TP_k$, $FN_k$ e $FP_k$ são, respectivamente, verdadeiros positivos,
falsos negativos e falsos positivos ao tomar $k$ como classe positiva.
A acurácia balanceada e a média macro de F1 foram calculadas conforme a
Equação~\eqref{eq:metricas_resumo_loso}:
\begin{equation}
    \mathrm{Acc}_{\mathrm{bal}}=
    \frac{\mathrm{Recall}_0+\mathrm{Recall}_1}{2},
    \qquad
    F1_{\mathrm{macro}}=\frac{F1_0+F1_1}{2}.
    \label{eq:metricas_resumo_loso}
\end{equation}
O recall da classe 0 descreve a identificação de leituras do grupo
inoculado; os denominadores das métricas são leituras, não vasos. As
leituras sucessivas de uma mesma coleta e as coletas repetidas do mesmo
vaso não constituem réplicas biológicas independentes. Em particular,
a acurácia reunida atribui mais peso a coletas com maior número de leituras.
